% 三版共用的代码块样式：Pygments 词法配色、minted 默认设置与代码卡片标题。
% 由各版布局在 shared/fonts.tex、shared/colors.tex 之后 \input。
% 某版需要不同的代码字号时，在 \input 之后用 \setminted{fontsize=…} 覆盖。

% 使用 bw 提供 Pygments 转义与分词接口，在每次加载样式后替换 token 外观。
% 兼容目标：minted 3.x / Pygments 2.x；不安装系统 Python 插件。
% 未列出的 token 不增加颜色、粗体、下划线或错误边框。
\makeatletter
\newcommand{\BookCodeBlueText}[1]{\textcolor{BookCodeBlue}{#1}}
\newcommand{\BookCodeGreenText}[1]{\textcolor{BookCodeGreen}{#1}}
\newcommand{\BookCodeCommentText}[1]{\textcolor{BookCodeComment}{#1}}
\def\BookToken@k{\let\PYG@tc\BookCodeBlueText}
\let\BookToken@m\BookToken@k
\let\BookToken@no\BookToken@k
\let\BookToken@ow\BookToken@k
% 内建名称：C++ 的 true、false、NULL，Bash、Python 的内建命令与函数。
\let\BookToken@nb\BookToken@k
\def\BookToken@s{\let\PYG@tc\BookCodeGreenText}
\def\BookToken@c{\let\PYG@tc\BookCodeCommentText}
\def\BookToken@cp{\let\PYG@tc\BookCodeBlueText\let\PYG@it\relax}
\def\BookToken@cpf{\let\PYG@tc\BookCodeGreenText\let\PYG@it\relax}
% Pygments 2.19、2.20 把 1ll、1llu、1lu 分成浮点数 1l 与名称 l、lu、u（pygments/pygments#1847），
% 把 C++23 的 1z、1uz、1zu 分成整数与名称 z、zu。数字后紧跟只由 l、u、z（含大写）组成的名称时，
% 名称只会是拆开的后缀，改按数字着色；10_km 等自定义后缀以下划线开头，不受影响。上游修复进入 TeX Live 后可删除。
% fvextra 给每个 token 加分组，折行时又以 \FancyVerbBreakStart/Stop 包住内容，故用全局状态，匹配时容许这层包装。
% 每行先测宽再排版，每遍都从 \FancyVerbFormatText 开始，在此清除上一遍留下的状态。
\ExplSyntaxOn
\bool_new:N \g__book_code_after_number_bool
\regex_const:Nn \c__book_code_suffix_regex
  { \A (?: \c{FancyVerbBreakStart} )? [lLuUzZ]{1,3} (?: \c{FancyVerbBreakStop} )? \Z }
\cs_new_protected:Npn \BookCode@Suffix #1#2
  {
    \bool_if:NT \g__book_code_after_number_bool
      {
        \str_if_eq:nnT {#1} {n}
          { \regex_match:NnT \c__book_code_suffix_regex {#2} { \BookToken@m } }
      }
    \str_if_eq:eeTF { \str_range:nnn {#1} { 1 } { 3 } } { l+m }
      { \bool_gset_true:N \g__book_code_after_number_bool }
      { \bool_gset_false:N \g__book_code_after_number_bool }
  }
\cs_new_eq:NN \__book_code_format_text:n \FancyVerbFormatText
\cs_set_protected:Npn \FancyVerbFormatText #1
  {
    \bool_gset_false:N \g__book_code_after_number_bool
    \__book_code_format_text:n {#1}
  }
\ExplSyntaxOff
% Pygments 的 Bash 词法把 <0.in、diff 1.out 中文件名开头的数字单独标为数字（Literal.Number），
% 其后的 .in、.out 是不经 \PYG 的普通文字；数字后紧跟“.”时改按普通文字排。
% fvextra 在每个 token 之后结束分组，故越过这个 \endgroup 再看下一个字符；不是 \endgroup 时照常着色。
\def\BookCode@Number#1{\@ifnextchar\endgroup{\BookCode@Number@i{#1}}{\PYG@do{#1}}}
\def\BookCode@Number@i#1\endgroup{%
  \@ifnextchar.{\let\PYG@tc\relax\PYG@do{#1}\endgroup}{\PYG@do{#1}\endgroup}}
\AddToHook{cmd/minted@defstyle/after}{%
  \edef\BookCodeStyle{\mintedtexoptvalueof{style}}%
  \ifdefstring{\BookCodeStyle}{bw}{%
    \def\PYG@tok#1{\ifcsname BookToken@#1\endcsname
      \csname BookToken@#1\endcsname\fi}%
    % 同 Pygments 的 \PYG，只加入上面两项处理。
    \def\PYG#1#2{\PYG@reset\PYG@toks#1+\relax+\BookCode@Suffix{#1}{#2}%
      \ifstrequal{#1}{l+m}{\BookCode@Number{#2}}{\PYG@do{#2}}}%
  }{}%
}
\makeatother
\setminted{
  fontfamily=bookcode,fontsize=\fontsize{8}{10.5}\selectfont,
  baselinestretch=1,style=bw,linenos=false,
  formatcom=\color{BookCodeText}\mdseries\upshape,
  breaklines=true,breakanywhere=false,tabsize=4
}
% 词法高亮并非 C++ 语义分析；仅补充明确要求的三个标准库类型名称、本书的整数类型别名 i64、i128、u64、u32，
% 以及 GCC 的 __int128_t、__uint128_t（__int128 已是关键字）。这些别名与 GCC 类型名不会用作变量名。
\setminted[cpp]{extrakeywordstype={vector,string,pair,i64,i128,u64,u32,__int128_t,__uint128_t}}
\setminted[c++]{extrakeywordstype={vector,string,pair,i64,i128,u64,u32,__int128_t,__uint128_t}}
\setminted[cxx]{extrakeywordstype={vector,string,pair,i64,i128,u64,u32,__int128_t,__uint128_t}}
\newfontfamily\BookCodeTitleLatin{PingFang SC Medium}
\newCJKfontfamily\BookCodeTitleCJK{PingFang SC Medium}
% coltext 同时保证 tcolorbox 续页恢复深灰，而非默认纯黑。
\tcbset{input-code-title/.style={coltext=BookCodeText,fonttitle=\mdseries\upshape\BookCodeTitleLatin\BookCodeTitleCJK\fontsize{9}{12}\selectfont}}
% 代码卡片与引用块共用的前后距，取代 tcolorbox 默认的 beforeafter skip balanced；
% 长度 \BookBoxBeforeSkip、\BookBoxAfterSkip 由各版布局设定。
% 盒前用 tcolorbox 的 before skip：盒顶距上一行的底为前距，盒与盒之间也是这个距离。
% 紧跟标题时 tcolorbox 不执行盒前设置（\iftcb@nobreak），标题与盒之间不受影响。
% 盒后把 \prevdepth 归零再加后距：后接正文时按段距与行距接排，后接标题时盒底相当于上一行的基线。
\makeatletter
\tcbset{book-box-spacing/.style={
  before skip=\BookBoxBeforeSkip,
  after={\tcb@parfillskip@check\par\ifvmode\prevdepth\z@\addvspace{\BookBoxAfterSkip}\fi}}}
\makeatother
